\documentclass[11pt,letterpaper]{article}
\usepackage[T1]{fontenc}
\usepackage{lmodern}
\usepackage[margin=1in,headheight=15pt]{geometry}
\usepackage{amsmath,amssymb,amsthm,mathtools}
\usepackage{microtype,booktabs,array,longtable}
\usepackage[dvipsnames]{xcolor}
\usepackage{enumitem}
\usepackage{fancyhdr}
\usepackage{listings}
\usepackage{hyperref}
\usepackage[nameinlink,noabbrev]{cleveref}
\definecolor{Ink}{HTML}{17324D}
\definecolor{Accent}{HTML}{176B78}
\definecolor{Pale}{HTML}{EDF4F6}
\hypersetup{colorlinks=true,linkcolor=Ink,citecolor=Accent,urlcolor=Accent,
 pdftitle={Sharp size thresholds for failures of log-concavity in independence systems},
 pdfauthor={ChatGPT, OpenAI},pdfsubject={Extremal refinements of the Xie--Xu augmentation conjecture}}
\pagestyle{fancy}
\fancyhf{}
\fancyhead[L]{\small\sffamily\color{Ink}Sharp log-concavity thresholds}
\fancyhead[R]{\small\sffamily Research note}
\fancyfoot[C]{\small\thepage}
\renewcommand{\headrulewidth}{0.3pt}
\setlength{\parskip}{4pt}
\setlength{\parindent}{1.1em}
\setlength{\emergencystretch}{2em}
\numberwithin{equation}{section}
\newtheorem{theorem}{Theorem}[section]
\newtheorem{proposition}[theorem]{Proposition}
\newtheorem{lemma}[theorem]{Lemma}
\newtheorem{corollary}[theorem]{Corollary}
\theoremstyle{definition}
\newtheorem{definition}[theorem]{Definition}
\newtheorem{problem}[theorem]{Problem}
\newtheorem{example}[theorem]{Example}
\theoremstyle{remark}
\newtheorem{remark}[theorem]{Remark}
\newcommand{\I}{\mathcal I}
\newcommand{\M}{\mathcal M}
\newcommand{\LC}{\mathrm{LC}}
\newcommand{\OLC}{\mathrm{OLC}}
\newcommand{\ULC}{\mathrm{ULC}}
\newcommand{\N}{\mathbb N}
\newcommand{\rank}{\operatorname{rank}}
\newcommand{\gap}{\operatorname{gap}}
\newcommand{\supp}{\operatorname{supp}}
\newcommand{\ind}{\operatorname{Ind}}
\DeclareMathOperator{\lcdef}{\Delta}
\lstset{basicstyle=\ttfamily\small,columns=fullflexible,keepspaces=true,
 breaklines=true,frame=single,rulecolor=\color{Pale},showstringspaces=false,
 keywordstyle=\color{Ink},commentstyle=\color{Accent},language=Python}

\begin{document}
\begin{titlepage}
\setlength{\parindent}{0pt}
\vspace*{0.65in}
{\sffamily\color{Accent}\large COMBINATORICS \quad / \quad EXACT EXTREMAL RESULTS\par}
\vspace{0.35in}
{\sffamily\bfseries\color{Ink}\LARGE Sharp size thresholds for failures\\[0.12em]
of log-concavity in\\[0.12em]
independence systems\par}
\vspace{0.2in}
{\Large A seven-element obstruction, rankwise optimality,\par
and a full-skeleton construction\par}
\vspace{0.4in}
{\large ChatGPT (OpenAI)\par}
{September 19, 2026\par}
\vspace{0.35in}
\begin{abstract}
A weak additive form of matroid augmentation does not force log-concavity.
The existence of counterexamples to Xie and Xu's rank-dependent conjecture
has already been announced by Alex Chengyu Li. We therefore address a
sharper extremal question: for each rank, how many ground-set elements are
necessary for failure? We prove that the unrestricted minima are attained
already by systems with exact augmentation parameter three. For ranks
$4,5,6$, the minimum sizes for ordinary log-concavity failure are $7,7,8$;
for every rank $r\ge7$ the minimum is $r+1$. Ordered log-concavity has minima
$6,7,r+1$ for $r=4$, $r=5$, and $r\ge6$, respectively. Ground-set-normalized
ultra log-concavity has minimum $r+1$ for every $r\ge4$.

The witnesses are a complete low-dimensional skeleton together with one
simplex. We classify all three concavity properties of this family by a
single coefficient inequality, establish global minimality using finite
Kruskal--Katona bounds, and obtain rational generating functions,
recurrences, and an exact vanishing log-concavity ratio. A standard-library
Python verifier checks the displayed examples, all possible face vectors
through seven active vertices, and more than two million uniform-family
shadow cases. The proofs are mathematical proofs in English, not a Lean
formalization. Historical priority for the refinements is not asserted.
\end{abstract}
\vfill
\noindent\colorbox{Pale}{\parbox{\dimexpr\textwidth-2\fboxsep\relax}{%
\textbf{Status and attribution.} This is an unrefereed research report.
It does \emph{not} claim the first disproof of Xie--Xu Conjecture~4.8.
The original conjecture, the recently located prior disproof, and the
sharp extremal question investigated here are separated explicitly in
Section~1. The archive contains the source, compiled article, exact
verification code, and generated certificates.}}
\end{titlepage}

\tableofcontents
\newpage

\section{The problem, its current status, and the main result}
\label{sec:problem}

\subsection{Independence systems and the augmentation convention}
An \emph{independence system} is a pair $\M=(E,\I)$, where $E$ is a finite
set and $\I\subseteq2^E$ contains $\varnothing$ and is downward closed:
if $T\in\I$ and $S\subseteq T$, then $S\in\I$.
Write $n=|E|$ and
\[
 r=\rank\M=\max\{|S|:S\in\I\},\qquad
 a_k=|\{S\in\I:|S|=k\}|.
\]
Our generating polynomial is $P_\M(x)=\sum_{k=0}^n a_kx^k$.
It has degree $r$, but we retain the zero coefficients through index $n$.
A vertex is \emph{active} when its singleton belongs to $\I$.

For a positive integer $l$, consider the following additive augmentation
condition:
\begin{equation}
\label{eq:augmentation}
 S,T\in\I,\quad |T|\ge |S|+l
 \quad\Longrightarrow\quad
 \text{some }x\in T\setminus S\text{ satisfies }S\cup\{x\}\in\I.
 \tag{$A_l$}
\end{equation}
Define $\lambda(\M)$ to be the least positive $l$ for which
\eqref{eq:augmentation} holds. Such an $l$ exists: $l=r+1$ makes the premise
impossible. The ordinary matroid augmentation axiom is $A_1$.
The phrase ``$l$-matroid'' below refers only to this additive convention,
not to other uses of similar terminology.

Xie and Xu introduce their terminology on printed page~14 of
\cite{XieXu}. Their displayed axiom $(*)$ writes $x\in T$ without explicitly
removing $S$. We require a genuinely new element, as in the usual
augmentation axiom. Every counterexample below satisfies this
\emph{stronger} interpretation; thus the omission cannot invalidate the
counterexamples. Exact values of $\lambda$ in this report always refer to
\eqref{eq:augmentation}.

\subsection{The three shapes, and a normalization that matters}
For a nonnegative sequence $(a_0,\ldots,a_n)$ we use the following
inequalities, at every $1\le k\le n-1$:
\begin{align}
\LC:\qquad &a_k^2\ge a_{k-1}a_{k+1},\label{eq:lc}\\
\OLC:\qquad &k a_k^2\ge (k+1)a_{k-1}a_{k+1},\label{eq:olc}\\
\ULC_n:\qquad &k(n-k)a_k^2
 \ge (k+1)(n-k+1)a_{k-1}a_{k+1}.\label{eq:ulc}
\end{align}
Ordered log-concavity means log-concavity of $(k!a_k)$.
Ultra log-concavity in \eqref{eq:ulc} means log-concavity of
$(a_k/\binom nk)$. These equivalences follow by substituting the factorials
and cancelling. Consequently $\ULC_n$ implies $\OLC$, which implies $\LC$.
The coefficient sequence of a nonempty independence system has no internal
zeros: a face of size $k$ supplies faces of every smaller size.

The subscript $n$ is essential. Normalizing by $\binom rk$, where $r$ is
only the degree, gives a different and stronger condition when $r<n$.
The sequence in Xie--Xu Conjecture~4.8 is explicitly indexed from $0$ to
$n$, so \eqref{eq:ulc} is the relevant normalization.

\subsection{A literature-status correction and the sharper target}
Xie--Xu Conjecture~4.8 asserts all three properties when
\begin{equation}
\label{eq:XX}
 \lambda(\M)\le \left\lfloor\frac{r}{2}\right\rfloor+1.
\end{equation}
The statement is on printed page~16 of the arXiv version \cite{XieXu}.
It was the initial target of this investigation.

A later status check located Li's September 2026 working-paper announcement
\cite{Li}. Its public overview already gives a rank-four example on
thirteen elements, with sequence $1,13,25,7,2$, and describes fixed-parameter
families with ratios tending to zero. Thus the existence question must
\emph{not} be presented here as still open or as first resolved by this
report. The full SSRN/Zenodo manuscript was not retrievable in this session;
its public overview and bibliographic record were available. Detailed
overlap with its unretrieved contents, and historical priority for the
refinements below, remain unverified.

The resulting research target is the following exact extremal problem.
It is a natural refinement, rather than a claim that the original conjecture
remained unsolved at the date of this report.

\begin{problem}\label{prob:sharp}
For each $r\ge4$, determine the least ground-set size at which a rank-$r$
independence system satisfying \eqref{eq:XX} can fail $\LC$, $\OLC$, or
$\ULC_n$. Can the minima already be attained with $\lambda=3$?
\end{problem}

For a property $Q\in\{\LC,\OLC,\ULC\}$, write $N_Q^{\rm all}(r)$ for the
least size of \emph{any} rank-$r$ independence system failing $Q$.
Write $N_Q^{\rm XX}(r)$ when \eqref{eq:XX} is imposed and
$N_Q^{(3)}(r)$ when $\lambda=3$ is imposed. In the ultra log-concave case,
the normalization is always the current ground-set size.

\begin{theorem}[Sharp rankwise thresholds]\label{thm:main}
For each $r\ge4$ and each of the three properties,
\[
 N_Q^{\rm all}(r)=N_Q^{\rm XX}(r)=N_Q^{(3)}(r).
\]
Their common values are as follows.
\begin{center}
\begin{tabular}{@{}c c c c@{}}
\toprule
Rank $r$ & Ordinary $\LC$ & Ordered $\OLC$ & Ultra $\ULC_n$\\
\midrule
$4$ & $7$ & $6$ & $5$\\
$5$ & $7$ & $7$ & $6$\\
$6$ & $8$ & $7$ & $7$\\
$r\ge7$ & $r+1$ & $r+1$ & $r+1$\\
\bottomrule
\end{tabular}
\end{center}
Every entry is attained by a system with all vertices active and with no
vertex belonging to every maximal independent set.
\end{theorem}

The upper bounds come from one explicit family in Section~3. The lower
bounds in Section~4 apply without an augmentation assumption. This is why
the unrestricted and restricted minima coincide. Globally, without fixing
the rank, the smallest possible ground-set sizes for failures of ordinary,
ordered, and ultra log-concavity are therefore $7,6,5$, respectively.

\section{The small obstruction, with a complete direct proof}
\label{sec:small}
Let $E=\{1,\ldots,7\}$ and $A=\{1,2,3,4\}$. Define
\begin{equation}\label{eq:seven}
 \I=2^A\ \cup\ \{S\subseteq E:|S|\le2\}.
\end{equation}
Thus every pair is independent, but the only independent sets of size
three or four are the subsets of $A$.

\begin{proposition}\label{prop:seven}
The system in \eqref{eq:seven} has rank $4$, exact augmentation parameter
$3$, and generating polynomial
\[
 P(x)=1+7x+21x^2+4x^3+x^4.
\]
It has $34$ independent sets and fails all three forms of log-concavity.
\end{proposition}
\begin{proof}
Both families in the union are downward closed. The set $A$ is independent
and no independent set has more than four elements, giving rank four.
Counting at each size yields the stated coefficients and $34$ total faces.

To verify $A_3$, take independent $S,T$ with $|T|\ge|S|+3$.
Since $|T|\le4$, we have $|S|\le1$. Any new element of $T\setminus S$
then produces a set of size at most two, which is independent. This proves
$A_3$ for every possible pair $S,T$.

To see that $A_2$ fails, use $S=\{1,5\}$ and $T=A$.
Their sizes differ by two. For every $x\in T\setminus S$, the three-set
$S\cup\{x\}$ contains $5$, so it is not independent. Hence $\lambda=3$.
This parameter satisfies $3=\lfloor4/2\rfloor+1$.

Finally, at index $3$ the ordinary inequality would require
$4^2\ge21\cdot1$, whereas $16<21$. The stronger inequalities fail as well.
\end{proof}

The same construction on five or six vertices separates the three
properties. It is useful to see their actual, integer-valued defects.
For a sequence on $n$ elements define
\begin{align*}
 D_{\LC}(k)&=a_k^2-a_{k-1}a_{k+1},\\
 D_{\OLC}(k)&=k a_k^2-(k+1)a_{k-1}a_{k+1},\\
 D_{\ULC}(k)&=k(n-k)a_k^2-(k+1)(n-k+1)a_{k-1}a_{k+1}.
\end{align*}
A negative defect is a certificate of failure.
\begin{table}[ht]
\centering
\begin{tabular}{@{}c l r r r@{}}
\toprule
$n$ & Nonzero coefficients & $D_{\LC}(3)$ & $D_{\OLC}(3)$ & $D_{\ULC}(3)$\\
\midrule
$5$ & $1,5,10,4,1$ & $6$ & $8$ & $-24$\\
$6$ & $1,6,15,4,1$ & $1$ & $-12$ & $-96$\\
$7$ & $1,7,21,4,1$ & $-5$ & $-36$ & $-228$\\
\bottomrule
\end{tabular}
\caption{Rank-four witnesses, all with exact augmentation parameter three.
All other tested indices satisfy the relevant inequalities.}
\label{tab:small}
\end{table}

For example, the five-element ultra inequality is
$3\cdot2\cdot16\ge4\cdot3\cdot10$, or $96\ge120$, which is false.
The six-element ordered inequality is $3\cdot16\ge4\cdot15$, or
$48\ge60$, also false.

\begin{remark}[Sparse variants]
The full pair skeleton is convenient, not required. From the five-element
example delete the pair $\{1,5\}$; from the six-element example delete
$\{1,5\},\{5,6\}$; from the seven-element example delete
$\{1,5\},\{5,6\},\{5,7\},\{6,7\}$.
The resulting sequences are
\[
 (1,5,9,4,1),\qquad (1,6,13,4,1),\qquad (1,7,17,4,1).
\]
They have $20,25,30$ faces and still have $\lambda=3$. Indeed, the only
nontrivial $A_3$ test beyond the empty set has a singleton $S$ and $T=A$;
every outside vertex still has a permitted pair with $A$. The same
$A_2$ obstruction can be chosen from a retained mixed pair. The last
sequence has the unit ordinary defect $16-17=-1$. We do not need a claim
of minimum total face count for the rankwise ground-size theorem.
\end{remark}

\section{The full-skeleton-plus-simplex family}
\label{sec:family}
Fix integers $n>r\ge2$ and $1\le s\le r-1$, with $A\subset E$, $|A|=r$,
$|E|=n$. Put
\begin{equation}\label{eq:family}
 \I_{n,r,s}=2^A\cup\{S\subseteq E:|S|\le s\}.
\end{equation}
The word ``skeleton'' will refer to the cutoff by \emph{cardinality} $s$;
in the usual topological dimension convention it is the
$(s-1)$-skeleton.

\begin{theorem}[Counts and exact augmentation]\label{thm:family}
The system \eqref{eq:family} has rank $r$, has no inactive vertices, and
satisfies
\begin{equation}\label{eq:coefficients}
 a_k=\begin{cases}
 \binom nk,&0\le k\le s,\\
 \binom rk,&s+1\le k\le r,\\
 0,&r<k\le n.
 \end{cases}
 \qquad
 \lambda(\I_{n,r,s})=r-s+1.
\end{equation}
Its maximal independent sets are $A$ and the $s$-sets not contained in $A$.
No vertex belongs to all maximal independent sets.
\end{theorem}
\begin{proof}
The downward-closure and rank statements follow directly from the
construction; counting gives \eqref{eq:coefficients}.
Let $l=r-s+1$. If $|T|\ge |S|+l$, then $|S|\le r-l=s-1$.
Thus adjoining any element of $T\setminus S$ produces a set of size at
most $s$. This verifies $A_l$.

Choose an $s$-set $S$ meeting $E\setminus A$ and put $T=A$.
The size difference is $r-s$. Every possible new element gives a mixed
$(s+1)$-set, absent from \eqref{eq:family}. Hence $A_{r-s}$ fails,
proving the exact parameter.

An independent set contained in $A$ extends to $A$. A mixed independent
set extends to a mixed $s$-set, and such an $s$-set cannot be extended.
This proves the description of maximal sets. Any outside vertex is omitted
by $A$. Given $v\in A$, choose a mixed $s$-set omitting $v$: there are at
least $r$ other available elements and $s\le r-1$. This maximal set omits
$v$, completing the proof.
\end{proof}

A vertex belongs to every maximal independent set if and only if it can be
adjoined to every independent set. To check the nontrivial direction,
extend an arbitrary face to a maximal face containing the vertex, and use
downward closure. Such a vertex is often called a cone vertex or a coloop
in this setting. Thus the theorem excludes a construction obtained merely
by adjoining universal independent elements.

When $s\ge2$, every pair is a face, so the one-skeleton is complete and the
simplicial complex is connected. For $s=r-2$ and $r\ge4$ it is not the
independence complex of a simple graph: all pairs are independent, but some
larger sets are not. Graph-realizable examples will be treated separately.

\subsection{Exactly one inequality can fail}
\begin{theorem}[Complete classification within the family]\label{thm:shapes}
For $1\le s\le r-2$, put
\[
 L=\binom ns,\qquad M=\binom r{s+1},\qquad H=\binom r{s+2}.
\]
The entire sequence \eqref{eq:coefficients} is
\begin{align}
 \LC &\quad\Longleftrightarrow\quad M^2\ge LH,\label{eq:testLC}\\
 \OLC &\quad\Longleftrightarrow\quad (s+1)M^2\ge(s+2)LH,\label{eq:testOLC}\\
 \ULC_n &\quad\Longleftrightarrow\quad
 (s+1)(n-s-1)M^2\ge(s+2)(n-s)LH.\label{eq:testULC}
\end{align}
In each case the only possible failing index is $s+1$.
For $s=r-1$, all three properties hold for every $n>r$.
\end{theorem}
\begin{proof}
For a positive sequence, log-concavity is equivalent to its successive
ratios being nonincreasing. In a binomial block, the three relevant ratio
sequences are obtained from
\[
 \frac{\binom n{k+1}}{\binom nk}=\frac{n-k}{k+1},\qquad
 \frac{(k+1)!\binom n{k+1}}{k!\binom nk}=n-k,\qquad
 \frac{\binom nk}{\binom nk}=1.
\]
They show all three inequalities strictly below the transition. For the
upper block, ordinary and ordered ratios are respectively
$(r-k)/(k+1)$ and $r-k$, again decreasing. The ultra-normalized upper block
has ratios
\begin{equation}\label{eq:upper-ratio}
 \frac{\binom r{k+1}/\binom n{k+1}}{\binom rk/\binom nk}
 =\frac{r-k}{n-k}.
\end{equation}
These decrease with $k$, since the difference between consecutive values
has numerator $n-r>0$.

At index $s$, the middle and previous terms are those of the lower binomial
block, while the next term has only decreased from $\binom n{s+1}$ to
$\binom r{s+1}$. The desired inequalities are therefore preserved.
Every index at least $s+2$ lies entirely in the upper block, or meets the
terminal zeros, and has already been covered. Substitution at index $s+1$
gives exactly \eqref{eq:testLC}--\eqref{eq:testULC}.

When $s=r-1$, the sole transition is a final downward change in $a_r$;
there is no following positive coefficient capable of making an inequality
at $s+1$ fail. The same block argument proves all three properties.
\end{proof}

\subsection{Parameter three and the explicit thresholds}
The decisive specialization is $s=r-2$, for $r\ge3$. It gives
\begin{equation}\label{eq:Pnr}
 P_{n,r}(x)=\sum_{k=0}^{r-2}\binom nk x^k+r x^{r-1}+x^r,
 \qquad \lambda=3.
\end{equation}
All three global tests reduce to the last nontrivial index $r-1$:
\begin{align}
 \LC &\Longleftrightarrow\binom n{r-2}\le r^2,\label{eq:3LC}\\
 \OLC &\Longleftrightarrow\binom n{r-2}\le r(r-1),\label{eq:3OLC}\\
 \ULC_n &\Longleftrightarrow
 \binom n{r-2}\le r(r-1)\frac{n-r+1}{n-r+2}.\label{eq:3ULC}
\end{align}
For every $r\ge4$, parameter three is allowed by \eqref{eq:XX}.
These are necessary and sufficient tests for the full sequence, not just
sufficient tests for failure.

The codimension-one case, $n=r+1$, is particularly economical.
Let $Q_r=P_{r+1,r}$. Then
\begin{equation}\label{eq:Qr}
 Q_r(x)=(1+x)^{r+1}-\binom r2 x^{r-1}-r x^r-x^{r+1}.
\end{equation}
Since $\binom{r+1}{r-2}=\binom{r+1}{3}$, cancellation in
\eqref{eq:3LC}--\eqref{eq:3ULC} gives the exact classification
\begin{equation}\label{eq:codimclass}
 \begin{array}{c|c}
 \text{Property of }Q_r\ (r\ge3)&\text{Ranks for which it holds}\\\hline
 \LC &3\le r\le6\\
 \OLC &3\le r\le5\\
 \ULC_{r+1}&r=3.
 \end{array}
\end{equation}
For ordinary log-concavity, for example, the test becomes
$r^2-6r-1\le0$, whose integral solutions with $r\ge3$ are precisely
$3,4,5,6$. For ordered log-concavity it becomes $r+1\le6$.
For ultra log-concavity it becomes $r+1\le4$.

\section{Global minimality, not just minimality within the construction}
\label{sec:minimality}
A counterexample does not by itself show that it is small in an absolute
sense. We now give finite, exact lower-bound certificates for all
independence systems needed to prove \cref{thm:main}.

\subsection{The shadow inequality}
For a family $\mathcal F$ of $q$-element sets, its \emph{lower shadow}
$\partial\mathcal F$ consists of all $(q-1)$-sets contained in at least
one member. If $m>0$, write its unique canonical $q$-binomial expansion
\begin{equation}\label{eq:KK-expansion}
 m=\binom{b_q}{q}+\binom{b_{q-1}}{q-1}+\cdots+\binom{b_j}{j},
 \qquad b_q>b_{q-1}>\cdots>b_j\ge j\ge1.
\end{equation}
Define
\begin{equation}\label{eq:shadow}
 \sigma_q(m)=\binom{b_q}{q-1}+\binom{b_{q-1}}{q-2}
 +\cdots+\binom{b_j}{j-1},\qquad \sigma_q(0)=0.
\end{equation}
The Kruskal--Katona shadow theorem gives
$|\partial\mathcal F|\ge\sigma_q(|\mathcal F|)$.
A primary-source account of the shadow theorem and its colex form appears
in the introduction to Bukh's paper \cite{Bukh}; the original references
are \cite{Kruskal,Katona}.

For an independence system, every shadow face is independent. Hence
\begin{equation}\label{eq:KK-needed}
 \sigma_{k+1}(a_{k+1})\le a_k.
\end{equation}
The expansion and its shadow can be evaluated by successive greedy choices
of the largest possible top index. For example,
\[
 26=\binom63+\binom42,\qquad
 \sigma_3(26)=\binom62+\binom41=19.
\]
This example is also one of the boundary values in the seven-vertex
enumeration in the archive.

\subsection{A small table that proves the lower bounds}
For a fixed $n$ and index $k$, let
\begin{equation}\label{eq:U}
 U_{n,k}(b)=\max\{c:0\le c\le\binom n{k+1},\quad
                         \sigma_{k+1}(c)\le b\}.
\end{equation}
If $a_k=b>0$, then
\begin{equation}\label{eq:ratio-bound}
 \frac{a_{k-1}a_{k+1}}{a_k^2}
 \le \frac{\binom n{k-1}U_{n,k}(b)}{b^2}.
\end{equation}
Multiply the right side by $1$, $(k+1)/k$, or
$(k+1)(n-k+1)/(k(n-k))$ for ordinary, ordered, or ultra log-concavity.
A maximum of at most one proves the corresponding inequality for
\emph{every} independence system under consideration.

For active ground sets, at $k=1$ the value of $b$ is fixed to $n$.
At $k\ge2$, maximize over $1\le b\le\binom nk$. Under the extra assumption
$\rank\M\ge t$, every $k$-subset of some $t$-face is independent; therefore
one may restrict to
\begin{equation}\label{eq:ranklower}
 \max\{1,\binom tk\}\le b\le\binom nk.
\end{equation}
\Cref{tab:bounds} records the resulting exact maxima. A dash denotes an
index outside the sequence.

\begin{table}[ht]
\centering\small
\setlength{\tabcolsep}{4.4pt}
\begin{tabular}{@{}l c c c c c c@{}}
\toprule
Case & $k=1$ & $k=2$ & $k=3$ & $k=4$ & $k=5$ & $k=6$\\
\midrule
$n=6$, $\LC$ & $5/12$ & $2/3$ & $15/16$ & $4/5$ & $5/12$ & ---\\
$n=5$, $\OLC$ & $4/5$ & $5/6$ & $5/6$ & $1/2$ & --- & ---\\
$n=4$, $\ULC_4$ & $1$ & $1$ & $1$ & --- & --- & ---\\
$n=7$, $r\ge6$, $\LC$ & $3/7$ & $28/45$ & $63/80$ & $14/15$ & $35/36$ & $3/7$\\
$n=6$, $r\ge5$, $\OLC$ & $5/6$ & $9/10$ & $1$ & $1$ & $1/2$ & ---\\
\bottomrule
\end{tabular}
\caption{Finite shadow bounds proving universal concavity in the indicated
cases. Every entry is at most one. For $\ULC_2$ and $\ULC_3$, the analogous
entries are all one. All maximizing pairs $(b,c)$ are in
\texttt{data/minimality\_bounds.csv}.}
\label{tab:bounds}
\end{table}

These are short finite arithmetic calculations, not probabilistic tests.
For example, the $15/16$ entry is the maximum of
$15U_{6,3}(b)/b^2$ over $1\le b\le20$, and is attained at $b=4$, $c=1$.
For $n=7$, $r\ge6$, and $k=5$, one has $b\ge\binom65=6$;
the maximum is $\binom74/6^2=35/36$ at $c=1$.
The complete algorithm defining every entry is \eqref{eq:KK-expansion}--
\eqref{eq:ranklower}, with no omitted search range. Appendix~A and the
verifier provide two ways to reproduce the finite calculations.

Zero middle coefficients cause no problem. Downward closure then implies
$a_{k+1}=0$, so all the relevant inequalities at $k$ hold.
Inactive ground elements can also be handled exactly. Removing them leaves
ordinary and ordered inequalities unchanged. For ultra log-concavity,
if the active size is $m\le n$, the factor
\[
 \frac{k+1}{k}\frac{n-k+1}{n-k}
\]
is no larger than the same factor with $m$ in place of $n$ whenever
$k<m$. Thus $\ULC_m$ on the active set implies $\ULC_n$ after adding
inactive elements. For active sizes $1,2,3,4$ the ultra bounds are vacuous
or the all-one rows described in \cref{tab:bounds}.

For ordinary and ordered lower bounds on ground sets smaller than the
first two rows, one can simply pad with inactive elements up to six or
five. The $k\ge2$ bounds remain valid because they use only the upper
ambient size; the index $1$ inequality follows directly from
$a_2\le\binom{a_1}{2}$. In the two rank-constrained rows, deleting an
inactive element leaves a full simplex when the rank equals the remaining
active size, which satisfies all the required inequalities.

\begin{proposition}[Universal small-size obstructions]\label{prop:lower}
Every independence system on at most six elements is ordinarily
log-concave. Every one on at most five elements is ordered log-concave.
Every one on at most four elements is ultra log-concave with its
actual ground-set normalization. Moreover, every rank-six system
on seven elements is ordinarily log-concave, and every rank-five system
on six elements is ordered log-concave.
\end{proposition}
\begin{proof}
Apply the at-most-one bounds in \cref{tab:bounds}, with the preceding
zero and inactive-element observations.
\end{proof}

\subsection{Completion of the sharp rankwise theorem}
\begin{proof}[Proof of \cref{thm:main}]
Any rank-$r$ system has $n\ge r$. If $n=r$, an independent $r$-set must
be the whole ground set, and downward closure forces the full simplex.
Its binomial sequence has all three properties. Therefore all desired
minima are at least $r+1$.

For ordinary log-concavity, \cref{prop:lower} increases this lower bound
to seven when $r=4$ or $5$, and to eight when $r=6$.
For ordered log-concavity, it increases the bound to six for $r=4$ and
to seven for $r=5$. No additional lower bound is needed for the other
entries.

For the matching upper bounds, use $\I_{n,r,r-2}$.
When $r=4$, choose $n=7,6,5$ for the three respective failures, as in
\cref{tab:small}. For $r=5$, the ordinary and ordered failures at $n=7$
follow from $\binom73=35>25>20$; ultra failure occurs at $n=6$ by
\eqref{eq:codimclass}. For $r=6$, ordinary failure at $n=8$ follows from
$\binom84=70>36$; the other two failures at $n=7$ again follow from
\eqref{eq:codimclass}. Every remaining entry is the codimension-one case
of that classification.

All these systems have $\lambda=3$, all vertices active, and no coloops
by \cref{thm:family}. For $r\ge4$, exact parameter three satisfies
\eqref{eq:XX}. The classes are nested, so their minima satisfy
\[
 N_Q^{\rm all}(r)\le N_Q^{\rm XX}(r)\le N_Q^{(3)}(r).
\]
The unrestricted lower bounds and parameter-three witnesses force equality
throughout.
\end{proof}

\section{Sequences, generating functions, and asymptotic strength}
\label{sec:sequences}

\subsection{The extremal size sequences}
The main theorem produces three elementary integer sequences, indexed by
$r\ge4$:
\[
\begin{array}{c|rrrrrrrr}
 r&4&5&6&7&8&9&10&11\\\hline
 N_{\LC}(r)&7&7&8&8&9&10&11&12\\
 N_{\OLC}(r)&6&7&7&8&9&10&11&12\\
 N_{\ULC}(r)&5&6&7&8&9&10&11&12.
\end{array}
\]
Their usefulness is their extremal interpretation; no claim of a new
numerical sequence or an existing OEIS accession is made.
Put $G_Q(z)=\sum_{j\ge0}N_Q(j+4)z^j$. Summing the eventual arithmetic
progression gives
\begin{align}
 G_{\ULC}(z)&=\frac{5-4z}{(1-z)^2},\label{eq:minGF}\\
 G_{\OLC}(z)&=\frac{5-4z}{(1-z)^2}+1+z,\\
 G_{\LC}(z)&=\frac{5-4z}{(1-z)^2}+2+z+z^2.
\end{align}
Thus the second differences vanish beyond the finitely many exceptional
ranks. These identities are formal power-series identities; analytically
they hold for $|z|<1$.

\subsection{Counting all faces in the codimension-one family}
Let $N_r=Q_r(1)$, where $Q_r$ is given by \eqref{eq:Qr}. For $r\ge3$,
\begin{equation}\label{eq:Nr}
 N_r=2^{r+1}-1-\binom{r+1}{2}.
\end{equation}
Indeed, the excluded faces are $\binom r2$ sets at size $r-1$, $r$ sets
at size $r$, and one set at size $r+1$. Their total is
$1+\binom{r+1}{2}$.

Extend the formula algebraically to $r=0,1,2$; no member of
\eqref{eq:family} with cutoff $r-2$ is being claimed at these indices.
The resulting sequence begins
\[
 1,2,4,9,21,48,106,227,475,978,1992,\ldots.
\]
Since $\sum_{r\ge0}\binom{r+1}{2}z^r=z/(1-z)^3$, we obtain
\begin{equation}\label{eq:totalGF}
 \sum_{r\ge0}N_rz^r
 =\frac2{1-2z}-\frac1{1-z}-\frac{z}{(1-z)^3}
 =\frac{1-3z+3z^2}{(1-2z)(1-z)^3}.
\end{equation}
Consequently
\begin{equation}\label{eq:recurrence}
 N_{r+4}=5N_{r+3}-9N_{r+2}+7N_{r+1}-2N_r\qquad(r\ge0).
\end{equation}
This follows either by multiplying \eqref{eq:totalGF} by its denominator
or by noting that $2^r$ and every polynomial of degree at most two are
annihilated by the characteristic polynomial
$(t-2)(t-1)^3$.

\subsection{A failure ratio tending to zero}
For the codimension-one sequence, the ratio at the failing position is
exactly
\begin{equation}\label{eq:ratio}
 R_r=\frac{a_{r-1}^2}{a_{r-2}a_r}
 =\frac{r^2}{\binom{r+1}{3}}=\frac{6r}{r^2-1}.
\end{equation}
Thus $R_r<1$ exactly when $r\ge7$, and $R_r\to0$.
More explicitly, for every $J\ge0$ and $r>1$,
\begin{equation}\label{eq:asymptotic}
 R_r=6\sum_{j=0}^{J}r^{-(2j+1)}
       +\frac{6r^{-(2J+3)}}{1-r^{-2}}.
\end{equation}
The remainder is positive and exact, so the asymptotic expansion
$6/r+6/r^3+6/r^5+\cdots$ needs no unquantified error estimate.
For $r\ge2$, the displayed remainder is at most
$8r^{-(2J+3)}$.

At the same time, $n=r+1$, $\lambda=3$, every singleton is independent,
and there are no coloops. In particular,
\[
 \frac{r}{n}\longrightarrow1,\qquad \frac{\lambda}{r}\longrightarrow0.
\]
The coefficient defect is the cubic integer expression
\begin{equation}\label{eq:defectcubic}
 a_{r-2}a_r-a_{r-1}^2
 =\frac{r^3-6r^2-r}{6},
\end{equation}
which starts $7,20,39,65,\ldots$ at $r=7,8,9,10,\ldots$.
Its fourth finite difference vanishes.

\begin{corollary}\label{cor:nouniversal}
No positive constant $c$ can make
$a_k^2\ge c\,a_{k-1}a_{k+1}$ hold uniformly for every independent-set
sequence with $\lambda=3$, even after requiring all vertices active,
no coloops, and $n=r+1$.
Nor can an upper bound on $\lambda$ depending only on rank restore
ordinary log-concavity if that bound permits $\lambda=3$ at arbitrarily
large ranks.
\end{corollary}
\begin{proof}
Use \eqref{eq:ratio} at the corresponding arbitrarily large ranks.
\end{proof}

These qualitative phenomena overlap with those already advertised in
Li's overview \cite{Li}. The contribution established here is the explicit
codimension-one construction, the sharp shape tests, and their use in the
exact ground-size theorem; this paragraph is not a priority claim for the
qualitative nonexistence of a uniform bound.

\section{What the augmentation parameter actually measures}
\label{sec:gap}
The following structural description is useful for understanding why a
rank-dependent hypothesis is too weak. It also explains the direct-sum
method mentioned in the prior-work overview.

For $U\subseteq E$, restrict the system to $U$ and let
$\mathcal B(U)$ be its inclusion-maximal independent sets. Define
\[
 w(U)=\max_{B\in\mathcal B(U)}|B|-\min_{B\in\mathcal B(U)}|B|.
\]
In particular $w(\varnothing)=0$.

\begin{proposition}[Hereditary maximal-set width]\label{prop:width}
For every finite independence system,
\begin{equation}\label{eq:width}
 \lambda(\M)=1+\max_{U\subseteq E}w(U).
\end{equation}
\end{proposition}
\begin{proof}
Let $W$ be the maximum on the right. Choose $U$ and maximal independent
sets $B,C$ of minimum and maximum size in $U$, so that $|C|-|B|=W$.
When $W\ge1$, no element of $C\setminus B$ can extend $B$, because $B$
is already maximal in $U$. Thus $A_W$ fails and $\lambda\ge W+1$.
For $W=0$, the same lower bound is simply $\lambda\ge1$.

Conversely, suppose independent $S,T$ admit no extension of $S$ by any
new element of $T$. In the restriction to $U=S\cup T$, the set $S$ is
maximal: otherwise a larger face containing it would supply a one-element
extension by downward closure. Extend $T$ to a maximal face $C$ in $U$.
Then
\[
 |T|-|S|\le |C|-|S|\le w(U)\le W.
\]
No unaugmentable pair can therefore have a size gap at least $W+1$.
This is $A_{W+1}$ and gives the reverse inequality.
\end{proof}

Thus $\lambda=3$ controls the spread of maximal-set sizes in every
restriction by two. It does not control the number of faces at adjacent
sizes. The family \eqref{eq:family} has an enormous collection of lower
faces and only the faces of one simplex above the cutoff; the width
condition alone cannot prevent that imbalance.

If systems on disjoint ground sets have independent families
$\I_1,\I_2$, their direct sum has independent sets $S_1\cup S_2$ with
$S_i\in\I_i$. Maximal faces in each restriction are formed independently
in the two components, so their widths add. Taking maxima in
\eqref{eq:width} yields the exact identity
\begin{equation}\label{eq:directsum}
 \lambda(\M_1\oplus\M_2)
 =\lambda(\M_1)+\lambda(\M_2)-1.
\end{equation}
Ranks add and generating polynomials multiply. In particular, a matroid
factor preserves the augmentation parameter of the other factor. This
identity explains how augmentation defects can remain fixed while rank
increases; it is not being advertised as a newly discovered method
relative to Li's prior overview.

\section{Graph independence polynomials and removal of isolated vertices}
\label{sec:graphs}
The principal minimal examples are not graph independence complexes.
It is nonetheless useful to verify that failure can occur in the graph
subclass too, without confusing independence and dependence polynomials.
For a graph $G$, let $I(G,x)$ count vertex sets containing no edge.

Let $H_m=K_m\vee\overline K_3$ be the join of a clique of size $m$ with
three mutually nonadjacent vertices: every vertex of the clique is joined
to each of those three vertices. A nonempty independent set is either a
subset of the three vertices or one clique vertex. Hence
\begin{equation}\label{eq:Hm}
 I(H_m,x)=(1+x)^3+mx.
\end{equation}
For $m\ge1$, its independence system has exact parameter three. Its rank
is three, a clique singleton cannot be extended toward the independent
triple, and $A_3$ has only the empty-set case to check.

\subsection{Arbitrarily large rank with a simple tail certificate}
Add $c=r-3$ isolated vertices. Their independence system is a free matroid,
so \eqref{eq:directsum} preserves parameter three. The new polynomial is
\begin{equation}\label{eq:graph-free}
 F_{r,m}(x)=(1+x)^r+mx(1+x)^{r-3}.
\end{equation}
Its top three coefficients are
\[
 a_{r-2}=\binom r2+m,\qquad a_{r-1}=r,\qquad a_r=1.
\]
The last ordinary log-concavity inequality fails precisely when
$m>\binom{r+1}{2}$. In particular, choosing
$m=\binom{r+1}{2}+1$ gives the tail $(r^2+1,r,1)$ and unit defect $-1$.
The graph has $r+m=\binom{r+2}{2}$ vertices. No minimum-graph-size claim is
made here.

\subsection{No isolated vertices and no cone vertices}
Instead of isolated vertices, take $c\ge1$ disjoint cliques of size
$t\ge2$. Each clique contributes a rank-one matroid with polynomial
$1+tx$. Thus
\begin{equation}\label{eq:graph-cliques}
 G_{c,m,t}=H_m\sqcup cK_t,
 \qquad I(G_{c,m,t},x)=((1+x)^3+mx)(1+tx)^c
\end{equation}
has rank $r=c+3$, exact parameter three, and no isolated graph vertices.
A graph vertex can be adjoined to every independent set exactly when it
is isolated, so these complexes also have no cone vertices.
The graph itself is generally disconnected; that is a different notion
from connectedness of the simplicial complex discussed earlier.

Reading the last three coefficients of \eqref{eq:graph-cliques} gives
\begin{align*}
 a_r&=t^c,\\
 a_{r-1}&=3t^c+ct^{c-1},\\
 a_{r-2}&=(3+m)t^c+3ct^{c-1}+\binom c2t^{c-2}.
\end{align*}
Since $t>0$, these expressions are valid also when $c=1$ (the last term
has zero coefficient). Direct expansion yields
\begin{equation}\label{eq:graph-defect}
 a_{r-1}^2-a_{r-2}a_r
 =t^{2c}\left(6-m+\frac{3c}{t}+\frac{c(c+1)}{2t^2}\right).
\end{equation}
For instance $c=1,m=8,t=2$ produces
$1,13,25,7,2$, the same sequence advertised in Li's prior overview,
with $7^2-25\cdot2=-1$.

For a uniform arbitrary-rank choice, set $m=7$ and $t=4c$. The factor in
parentheses is
\[
 -1+\frac34+\frac{c+1}{32c}\le-\frac3{16}<0.
\]
This gives rank-$r$ examples with no isolated vertices on
$10+4(r-3)^2$ vertices for every $r\ge4$.
These formulas provide independent, elementary verification of the graph
variants; they do not assert priority over the prior direct-sum method.

\section{A real-rootedness obstruction and the two normalizations}
\label{sec:roots}
Even a member of \eqref{eq:Pnr} that satisfies $\ULC_n$ need not be
real-rooted. There is a short proof requiring only differentiation and
Rolle's theorem.

\begin{proposition}\label{prop:roots}
For every $n>r\ge3$, the polynomial $P_{n,r}$ has at least one nonreal
root.
\end{proposition}
\begin{proof}
If all roots were real, repeated differentiation would preserve that
property. The derivative of order $r-2$ is the quadratic
\[
 P_{n,r}^{(r-2)}(x)
 =(r-2)!\binom n{r-2}+r(r-1)!x+\frac{r!}{2}x^2.
\]
Its discriminant is
\[
 r(r-1)(r-2)!^2
 \left(r(r-1)-2\binom n{r-2}\right).
\]
A nonnegative discriminant would require
$\binom n{r-2}\le r(r-1)/2=\binom r{r-2}$, which is impossible because
$n>r$ and $r-2\ge1$.
\end{proof}

For example, the $r=3,n=4$ polynomial $1+4x+3x^2+x^3$ satisfies
$\ULC_4$ by \eqref{eq:codimclass}, yet it has nonreal roots by the
proposition. This does not conflict with any implication from
real-rootedness to normalized coefficient inequalities. Normalization by
the degree $r$ is stronger than normalization by the ambient $n$:
at the last index of \eqref{eq:Pnr}, the degree-normalized ultra inequality
would require the very bound $\binom n{r-2}\le\binom r2$ just ruled out.

\section{Exact verification, reproducibility, and remaining scope}
\label{sec:verification}

\subsection{What was actually executed}
The archive's \texttt{verify.py} uses only the Python standard library.
Its full run completed successfully with Python 3.13.5. All comparisons
use integers or exact rational numbers; no numerical root approximation
or floating-point sign decision supports a theorem.

The verifier produced the counts below; their JSON summary is included
in the \texttt{data/} directory:
\begin{center}
\begin{tabular}{@{}l r@{}}
\toprule
Check & Cases\\
\midrule
General family shape classifications, $n\le24$ & $2{,}024$\\
Explicit face-system and augmentation checks, $n\le8$ & $56$\\
Direct-sum parameter checks & $6$\\
Possible face vectors on $1$ through $7$ active vertices & $5{,}459$\\
Finite maximal-ratio certificates & $26$\\
Graph polynomial identities & $54$\\
Unrestricted uniform-family shadow cases, $n\le6$ & $1{,}116{,}354$\\
Core-constrained uniform-family cases & $1{,}116{,}352$\\
\bottomrule
\end{tabular}
\end{center}
The last two rows total $2{,}232{,}706$ uniform-family cases.
These are not claimed to be that many distinct independence systems.

The core-constrained check independently examines all uniform families
containing the faces of a fixed six-element core on seven vertices, and
of a fixed five-element core on six vertices, at the levels needed by the
two exceptional rank lower bounds. It avoids enumerating all labelled
complexes on seven vertices. The largest optional uniform layer has
$20$ members, so this independent shadow check has at most $2^{20}$ states
in a single layer.

The separate face-vector enumeration uses the Kruskal--Katona criterion.
For each possible $a_k$, it recursively accepts exactly the values with
$\sigma_k(a_k)\le a_{k-1}$. These inequalities are sufficient as well as
necessary: choose the initial colex segment at each level; its shadow is
an initial colex segment of the preceding level, and the inequalities
ensure containment. The enumeration therefore concerns possible
\emph{vectors}, not just hypothetical arrays, while distinct complexes
with the same vector are counted once.

\begin{center}
\begin{tabular}{@{}r r r r r@{}}
\toprule
Active vertices & Possible vectors & Fail $\LC$ & Fail $\OLC$ & Fail $\ULC_n$\\
\midrule
1&1&0&0&0\\
2&2&0&0&0\\
3&5&0&0&0\\
4&16&0&0&0\\
5&70&0&0&5\\
6&457&0&3&47\\
7&4,908&23&120&809\\
\bottomrule
\end{tabular}
\end{center}
These counts are useful checks but are not substituted for the stated
rankwise theorem. In particular, counts of shape failures do not by
themselves assert that all failing systems obey $A_3$.

\subsection{Reproduction}
From the archive root, run
\begin{lstlisting}[language={}]
python verify.py --full --output regenerated-data
pdflatex -interaction=nonstopmode -halt-on-error article.tex
pdflatex -interaction=nonstopmode -halt-on-error article.tex
\end{lstlisting}
The first command regenerates the mathematical certificates. It records
the local Python version, so that metadata field may differ across
machines while the mathematical data agree. The LaTeX file is
self-contained, including its bibliography, and needs no BibTeX run.
The accompanying \texttt{README.md} describes the files, bit-mask
convention, and expected output.

\subsection{What is and is not settled}
The exact extremal problem \cref{prob:sharp} is resolved by
\cref{thm:main}, with explicit witnesses and universal lower bounds.
The three-parametric family's concavity classification and all asymptotic
claims have proofs independent of the computations. The finite lower
bounds use the classical shadow theorem plus the completely specified
finite arithmetic table; the exhaustive uniform-family checks give a
second verification for the specific sizes needed.

This does not prove or disprove the parameter-two conjecture: every main
counterexample has \emph{exact} parameter three. Xie and Xu separately
formulate the parameter-two assertion as Conjecture~4.4 in \cite{XieXu}.
The present construction with $s=r-1$ has parameter two and satisfies all
three shapes, but it is only a subclass of parameter-two systems.
Nor does an example about a graph's \emph{independent} sets refute the
same paper's Conjecture~4.1 about its \emph{dependent} sets. These are
separate counting problems.

Finally, the original existence disproof is prior work, not a new result
claimed here. The search was not a complete priority audit, and the
unretrieved full Li manuscript leaves possible overlap in refinements
unresolved. The strongest unconditional output of this report is the
proved sharp theorem and its reproducible certificates, not a claim of
publication priority or independent peer review.

\appendix
\section{Finite arithmetic certificates and independent shadow checks}
\label{app:certificates}

\subsection{The entire maximal-ratio calculation}
The following pseudocode is the finite proof behind each entry of
\cref{tab:bounds}. The function \texttt{shadow} is the explicit
binomial expansion \eqref{eq:KK-expansion} followed by
\eqref{eq:shadow}; it does not call an optimization oracle.
\begin{lstlisting}
for k in range(1, n):
    allowed_b = [n] if k == 1 else range(max(1, binom(t, k)),
                                        binom(n, k) + 1)
    maximum = 0
    for b in allowed_b:
        c = max(c for c in range(binom(n, k + 1) + 1)
                if shadow(c, k + 1) <= b)
        value = factor(kind, n, k) * binom(n, k - 1) * c / b**2
        maximum = max(maximum, value)
    assert maximum <= 1
\end{lstlisting}
Here $t=0$ means no rank lower bound. The actual program uses
\texttt{Fraction}, and its checking function remains active under
\texttt{python -O}. The displayed divisions in the pseudocode must not
be implemented with floating-point arithmetic when an exact certificate
is desired.

For illustration, the complete shadow minima for $q=4$ on six vertices,
indexed by the number $m=0,1,\ldots,15$ of four-sets, are
\[
 0,4,7,9,10,10,13,15,16,16,18,19,19,20,20,20.
\]
Thus one four-set forces four triples, five four-sets force ten triples,
and six four-sets force thirteen triples. Such small discontinuities are
why the exact shadow bound is substantially sharper than merely counting
incidences between adjacent levels.

\subsection{An independent enumeration of lower shadows}
Represent every subset of $\{1,\ldots,n\}$ by an integer mask, with
vertex $j$ represented by bit $j-1$. First enumerate all $q$-subsets and
all $(q-1)$-subsets. A $q$-set then has its own lower-shadow bit mask in
the $(q-1)$-set universe. For a family mask $F\ne0$, remove its lowest set
bit $b$ and use
\[
 \operatorname{shadowMask}(F)
 =\operatorname{shadowMask}(F\setminus\{b\})
       \mathbin{\mathrm{OR}}\operatorname{shadowMask}(\{b\}).
\]
Group the resulting popcounts by $|F|$ and take minima. This is independent
of the canonical binomial expansion. For all unrestricted layers on at
most six vertices, the computed minima agree exactly with $\sigma_q(m)$.

For a fixed core, initialize the shadow mask with all its required lower
faces and enumerate only the optional upper faces. Taking the minimum
lower count for each optional upper count yields a direct finite check of
the two exceptional rank rows. Every rank-$r$ system on $r+1$ vertices
contains an $r$-face; relabel it as the prescribed core. Therefore the
core enumeration covers every system needed for those two rows, even
though it does not impose all compatibility conditions between optional
layers. Omitting those extra conditions only enlarges the tested class,
which is safe for a universal lower bound.

\section{Compact classification of the smallest failing face vectors}
The following tables record all failing vectors at the three first
possible sizes, before any augmentation restriction is imposed. Terminal
zeros through the ambient size are suppressed.

At five vertices, the ultra log-concavity failures are
\[
 (1,5,3,1),\quad(1,5,6,4),\quad(1,5,6,4,1),\quad
 (1,5,9,4,1),\quad(1,5,10,4,1).
\]
At six vertices, the ordered failures are exactly
\[
 (1,6,b,4,1),\qquad b=13,14,15.
\]
The full list of twenty-three ordinary failures on seven vertices can be
compressed as follows:
\begin{center}
\begin{tabular}{@{}l l r@{}}
\toprule
Vector pattern & Parameter range & Number\\
\midrule
$(1,7,b,4,1)$ & $17\le b\le21$ & $5$\\
$(1,7,21,10,5)$ & --- & $1$\\
$(1,7,21,10,5,1)$ & --- & $1$\\
$(1,7,19,26,5,1)$ & --- & $1$\\
$(1,7,20,c,5,1)$ & $26\le c\le30$ & $5$\\
$(1,7,21,c,5,1)$ & $26\le c\le35$ & $10$\\
\bottomrule
\end{tabular}
\end{center}
The list immediately provides another check that none of the seven-vertex
ordinary failures has rank six, and none of the six-vertex ordered
failures has rank five. These lists are consequences of the finite
face-vector criterion, not a classification of labelled systems or their
augmentation parameters.

\begin{thebibliography}{9}
\bibitem{XieXu}
Yan-Ting Xie and Shou-Jun Xu,
\emph{Ultra log-concavity and real-rootedness of dependence polynomials},
arXiv:2408.09152v1, August 17, 2024.
\href{https://arxiv.org/abs/2408.09152}{arXiv record}.
The relevant locations are the sequence definitions on printed page~3,
the additive augmentation condition on page~14, and Conjecture~4.8 on
page~16. The record was checked on September 19, 2026.

\bibitem{Li}
Alex Chengyu Li,
\emph{Additive Augmentation Gaps Do Not Force Log-Concavity},
working paper, 2026.
\href{https://doi.org/10.2139/ssrn.7434141}{DOI: 10.2139/ssrn.7434141}.
Author's public overview:
\href{https://crabresearch.com/research/augmentation-gap-log-concavity?lang=en}{Crab Research, augmentation-gap-log-concavity}.
That page reports first public release September 6, 2026 and revision
September 18, 2026. Overview accessed September 19, 2026; the full linked
manuscript was not retrieved in this investigation.

\bibitem{Bukh}
Boris Bukh,
\emph{Multidimensional Kruskal--Katona theorem},
arXiv:1009.2375v2, 2011.
\href{https://arxiv.org/abs/1009.2375}{arXiv record}.
The introduction states the classical colex shadow theorem and the
Lov\'asz form, and gives the original references below.

\bibitem{Kruskal}
Joseph B. Kruskal,
\emph{The number of simplices in a complex},
in \emph{Mathematical Optimization Techniques},
University of California Press, Berkeley, 1963, pp.~251--278.

\bibitem{Katona}
Gyula Katona,
\emph{A theorem of finite sets},
in \emph{Theory of Graphs} (Proceedings of the Colloquium, Tihany, 1966),
Academic Press, New York, 1968, pp.~187--207.
\end{thebibliography}
\end{document}
